\chapter{Implementation}
\label{ch:implementation}

% ============================================================================
% This chapter describes the implementation: architecture, caching, pipeline
% ============================================================================

\section{Overview}
\label{sec:impl_overview}

This chapter describes the implementation of the sleep stage classification pipeline. The system is implemented in Python with a modular architecture consisting of 25 modules totaling approximately 16,571 lines of code.

\section{System Architecture}
\label{sec:architecture}

The pipeline is organized into the following modules:

\begin{itemize}
    \item \textbf{Data Loading:} \texttt{data\_loader\_boas.py} (703 lines)
    \item \textbf{Preprocessing:} \texttt{preprocessing.py} (431 lines)
    \item \textbf{Feature Extraction:} \texttt{feature\_extractor.py} (753 lines)
    \item \textbf{Feature Selection:} \texttt{feature\_selection.py} (942 lines)
    \item \textbf{Models:} \texttt{models.py} (535 lines)
    \item \textbf{Training:} \texttt{training.py} (1,226 lines)
    \item \textbf{Cross-Validation:} \texttt{cross\_validation.py} (506 lines)
    \item \textbf{Evaluation:} \texttt{evaluation.py} (601 lines)
    \item \textbf{Caching:} \texttt{feature\_cache.py} (304 lines), \texttt{loso\_cache.py} (598 lines)
    \item \textbf{Fingerprinting:} \texttt{fingerprint.py} (423 lines)
    \item \textbf{Utilities:} \texttt{utils.py} (432 lines), \texttt{config.py} (348 lines)
    \item \textbf{Visualization:} \texttt{visualization.py} (843 lines), \texttt{output\_formatter.py} (656 lines)
\end{itemize}

The entry points are:
\begin{itemize}
    \item \texttt{run\_experiment.py} (368 lines): Data loading and feature extraction (Stages 1--2)
    \item \texttt{run\_training.py} (752 lines): Model training and evaluation (Stage 3)
    \item \texttt{interactive\_menu.py} (571 lines): Interactive configuration interface
\end{itemize}

\section{Two-Tier Caching System}
\label{sec:caching}

A key contribution of this implementation is the two-tier caching system that dramatically reduces computational overhead for iterative experimentation.

\subsection{Motivation}
\label{subsec:cache_motivation}

Sleep stage classification experiments involve computationally expensive operations:

\begin{itemize}
    \item \textbf{Feature extraction:} Processing 128 subjects $\times$ $\sim$936 epochs $\times$ 6 channels $\times$ 30 seconds
    \item \textbf{LOSO training:} Training 128 models per configuration
\end{itemize}

Without caching, each parameter change requires recomputing all features and retraining all models. The two-tier cache addresses this inefficiency.

\subsection{Tier 1: Feature Cache}
\label{subsec:feature_cache}

The feature cache stores extracted features for each subject with a fingerprint based on:

\begin{itemize}
    \item EEG channel configuration (which channels to use)
    \item Preprocessing parameters (filter cutoffs, epoch duration)
    \item Feature extraction parameters (Welch window size, frequency bands)
    \item Code version (fingerprint of \texttt{feature\_extractor.py})
\end{itemize}

\textbf{Cache Hit Criteria:} If all parameters match and the subject's raw data has not changed, cached features are loaded instead of recomputed.

\textbf{Granularity:} Subject-level (128 independent cache entries)

\textbf{Storage:} \texttt{results/features\_cache\_global/subject\_XXX/}

\subsection{Tier 2: LOSO Model Cache}
\label{subsec:loso_cache}

The LOSO model cache stores trained models for each fold (held-out subject) with a fingerprint based on:

\begin{itemize}
    \item Random seed
    \item Code version
    \item Model type (XGBoost, Random Forest)
    \item Model hyperparameters
    \item Feature selection configuration (correlation threshold, top-K)
    \item Held-out subject ID
\end{itemize}

\textbf{Cache Hit Criteria:} If all parameters match and the training data (from feature cache) has not changed, the cached model is loaded instead of retrained.

\textbf{Granularity:} Per-fold (128 folds per configuration)

\textbf{Storage:} \texttt{results/loso\_cache/}

\subsection{Fingerprint Generation}
\label{subsec:fingerprints}

Fingerprints use SHA-256 cryptographic hashing to uniquely identify configurations:

\begin{lstlisting}[language=Python]
import hashlib

def generate_fingerprint(config_dict):
    """Generate SHA-256 fingerprint from configuration."""
    config_str = str(sorted(config_dict.items()))
    return hashlib.sha256(config_str.encode()).hexdigest()[:16]
\end{lstlisting}

This ensures:
\begin{itemize}
    \item \textbf{Collision resistance:} Different configs $\Rightarrow$ different fingerprints
    \item \textbf{Determinism:} Same config $\Rightarrow$ same fingerprint
    \item \textbf{Cache invalidation:} Any parameter change invalidates the cache
\end{itemize}

\subsection{Cache Efficiency}
\label{subsec:cache_efficiency}

% TODO: Add table showing cache hit rates during development
% Example: After initial run, 100% feature cache hits, 0% model cache hits
% After parameter change, 100% feature cache hits, variable model cache hits

The two-tier design allows:
\begin{itemize}
    \item Changing feature selection $\Rightarrow$ reuse features, retrain models
    \item Changing model hyperparameters $\Rightarrow$ reuse features, retrain models
    \item Unchanged configuration $\Rightarrow$ reuse everything (near-instant results)
\end{itemize}

\section{Pipeline Implementation}
\label{sec:pipeline_impl}

\subsection{Stage 1: Data Loading}
\label{subsec:data_loading}

The \texttt{data\_loader\_boas.py} module loads raw EEG data from the BOAS dataset. It handles:

\begin{itemize}
    \item Reading EDF/HDF5/SNIRF files
    \item Extracting specified EEG channels
    \item Loading sleep stage annotations
    \item Validating data integrity
\end{itemize}

\subsection{Stage 2: Feature Extraction}
\label{subsec:feature_extraction_impl}

The \texttt{feature\_extractor.py} module extracts features with parallel processing:

\begin{lstlisting}[language=Python]
from joblib import Parallel, delayed

features = Parallel(n_jobs=n_workers, backend='multiprocessing')(
    delayed(extract_epoch_features)(epoch_data, channels)
    for epoch_data in epochs
)
\end{lstlisting}

\textbf{Parallelization:} Epoch-level parallelization using \texttt{joblib.Parallel} with multiprocessing backend.

\textbf{Batch Size:} Configurable (default: 100 epochs/batch) to balance memory and CPU efficiency.

% TODO: Add feature extraction performance metrics
% Time per epoch, speedup from parallelization

\subsection{Stage 3: Training Pipeline}
\label{subsec:training_impl}

The \texttt{training.py} module orchestrates LOSO cross-validation:

\begin{algorithm}[htbp]
\caption{LOSO Cross-Validation with Model Cache}
\label{alg:loso}
\begin{algorithmic}[1]
\FOR{each subject $s$ in $1 \ldots 128$}
    \STATE $\text{fingerprint} \leftarrow \text{generate\_fingerprint}(\text{config}, s)$
    \IF{cache contains fingerprint}
        \STATE Load cached model and predictions
    \ELSE
        \STATE $X_{\text{train}}, y_{\text{train}} \leftarrow$ data from all subjects except $s$
        \STATE $X_{\text{test}}, y_{\text{test}} \leftarrow$ data from subject $s$
        \STATE Apply feature selection on $X_{\text{train}}$
        \STATE Train model on $X_{\text{train}}, y_{\text{train}}$
        \STATE Predict on $X_{\text{test}}$
        \STATE Cache model and predictions
    \ENDIF
    \STATE Evaluate predictions
\ENDFOR
\STATE Aggregate results across all 128 folds
\end{algorithmic}
\end{algorithm}

\section{Feature Selection Implementation}
\label{sec:feature_selection_impl}

The \texttt{feature\_selection.py} module implements the two-stage selection pipeline:

\begin{enumerate}
    \item \textbf{Correlation filtering:} Remove features with $r > \rho$
    \item \textbf{ANOVA selection:} Rank by F-score, select top-$k$
\end{enumerate}

\textbf{Benchmark:} ANOVA is 200$\times$ faster than mutual information with <1\% accuracy loss (see \texttt{benchmarks\_and\_tests/benchmark\_anova\_vs\_mi.py}).

\section{Model Implementation}
\label{sec:model_impl}

\subsection{XGBoost}
\label{subsec:xgboost_impl}

XGBoost is implemented using the \texttt{xgboost} Python library:

\begin{lstlisting}[language=Python]
import xgboost as xgb

model = xgb.XGBClassifier(
    max_depth=6,
    learning_rate=0.1,
    n_estimators=100,
    random_state=42
)
model.fit(X_train, y_train)
\end{lstlisting}

\subsection{Random Forest}
\label{subsec:rf_impl}

Random Forest is implemented using \texttt{scikit-learn}:

\begin{lstlisting}[language=Python]
from sklearn.ensemble import RandomForestClassifier

model = RandomForestClassifier(
    n_estimators=100,
    max_depth=20,
    random_state=42,
    n_jobs=-1
)
model.fit(X_train, y_train)
\end{lstlisting}

\section{Evaluation Implementation}
\label{sec:eval_impl}

The \texttt{evaluation.py} module computes metrics using \texttt{scikit-learn}:

\begin{itemize}
    \item Accuracy: \texttt{sklearn.metrics.accuracy\_score}
    \item Cohen's Kappa: \texttt{sklearn.metrics.cohen\_kappa\_score}
    \item F1-Score: \texttt{sklearn.metrics.f1\_score} (macro-averaged)
    \item Confusion Matrix: \texttt{sklearn.metrics.confusion\_matrix}
\end{itemize}

Results are aggregated across all 128 LOSO folds using mean and standard deviation.

\section{Visualization and Reporting}
\label{sec:visualization_impl}

The \texttt{visualization.py} module generates:

\begin{itemize}
    \item Confusion matrices (per-configuration and aggregated)
    \item Feature importance plots
    \item Performance comparison bar charts
    \item Cache hit rate visualizations
\end{itemize}

The \texttt{output\_formatter.py} module provides a formatted terminal UI with:

\begin{itemize}
    \item Configuration summary boxes
    \item Real-time fold progress with cache status icons (💾 = cache hit, ⏳ = cache miss)
    \item Color-coded performance metrics
\end{itemize}

\section{Configuration Management}
\label{sec:config_impl}

The \texttt{config.py} module provides a dataclass-based configuration system:

\begin{lstlisting}[language=Python]
from dataclasses import dataclass

@dataclass
class ExperimentConfig:
    model_type: str
    random_state: int
    feature_selection: FeatureSelectionConfig
    model_params: dict
\end{lstlisting}

Configurations are serialized to JSON for reproducibility.

\section{Dependencies}
\label{sec:dependencies}

Key Python libraries:

\begin{itemize}
    \item \textbf{Data processing:} numpy, pandas, scipy
    \item \textbf{Machine learning:} scikit-learn, xgboost
    \item \textbf{Signal processing:} scipy.signal
    \item \textbf{Parallelization:} joblib
    \item \textbf{Visualization:} matplotlib, seaborn
    \item \textbf{Utilities:} tqdm, logging
\end{itemize}

% TODO: Add requirements.txt content or version table

\section{Testing}
\label{sec:testing}

The implementation includes comprehensive tests:

\begin{itemize}
    \item \texttt{test\_loso\_cache.py} (1,055 lines): LOSO cache validation
    \item \texttt{test\_cache\_comprehensive.py} (292 lines): Feature cache tests
    \item \texttt{test\_loso\_cache\_fixes.py} (341 lines): Cache invalidation tests
    \item \texttt{test\_validation.py} (27 lines): Cross-validation tests
\end{itemize}

Benchmarks:

\begin{itemize}
    \item \texttt{benchmark\_anova\_vs\_mi.py} (245 lines): ANOVA vs mutual information
    \item \texttt{benchmark\_feature\_selection.py} (624 lines): Feature selection performance
\end{itemize}

\section{Summary}
\label{sec:impl_summary}

This chapter described the implementation of the sleep stage classification pipeline, with emphasis on the two-tier caching system using SHA-256 fingerprints for cache invalidation. The modular architecture with 25 Python modules totaling 16,571 lines provides a complete, reproducible system for sleep stage classification with computational efficiency optimizations.
